\section{Introduction}
\label{sec:intro}

Large language models have become working infrastructure of public
administration within a few years. An OECD survey of core government functions
documents generative-AI adoption across policymaking and service delivery,
including assistants that help public-service advisers retrieve information and
the sources behind it~\citep{oecd2024governing}. Air pollution policy is a domain
where this carries direct stakes for public health. Fine particulate matter is
among the leading environmental risk factors for premature mortality
worldwide~\citep{gbd2019risk,who2021aqg}, and the burden falls disproportionately
on the Global South. When a public environmental manager asks a model about an
ambient air quality standard, an attainment deadline, or the evidence behind an
intervention, an inaccurate answer can propagate into a regulatory brief, a
procurement decision, or an emergency response during a pollution episode.

This paper asks whether models answer such questions as reliably for Global
South countries as for Global North ones, and whether the way a query is framed
changes the answer. Both premises have support. Geographic factual bias is
measurable: error rates are $1.5\times$ higher for Sub-Saharan African than for
North American countries across 20 models~\citep{moayeri2024worldbench}, model
ratings on subjective topics track socioeconomic conditions up to
$\rho=0.70$~\citep{manvi2024llm}, and models systematically privilege statements
about the Global North~\citep{mirza2024global}. The linguistic asymmetry behind
it is structural: the languages of most Global South countries occupy the
under-served classes of the resource hierarchy~\citep{joshi2020state}, and large
Internet-derived corpora ``overrepresent hegemonic viewpoints and encode biases
potentially damaging to marginalized populations''~\citep{bender2021stochastic}.
Decolonial scholarship reads that asymmetry as an uneven epistemic
infrastructure reproduced by AI systems~\citep{mohamed2020decolonial}.

Three gaps organise the study. The first is narrower than it is usually
stated. Regulatory information with official ground truth has been benchmarked
once, by \citet{dahl2024legal}, against United States federal case law, and that
work is the closest precedent to ours. What has not been done is the
cross-national version: a register-anchored benchmark that is stratified for
equity across countries rather than run inside one jurisdiction, in a domain
whose keys are national instruments. The other evaluations that do span
countries use World Bank recall, subjective ratings, journalistic
fact-checking, or everyday cultural knowledge, none of which measures the tasks
a manager delegates, such as stating a binding standard, retrieving an
officially measured concentration, summarising health-burden evidence,
identifying legal instruments and executing agencies, or recommending an
operational response. Persona framing is unexamined as an
experimental lever, although practitioners routinely prefix a role and the
vendors advise it (Section~\ref{sec:related}); whether it reduces the gap by
steering toward regulation-compliant answers or amplifies it by inviting
confident fabrication has not been tested in a pre-specified design. And the
corpus-representation mechanism is usually asserted rather than measured against
confounders, and rarely separated into its two channels, the size of the
language corpus~\citep{xue2021mt5,abadji2022oscar} and the breadth of
country coverage, a distinction that matters here because the geographic gap is
measured in English.

\subsection*{Objectives and scope}

The study has three bounded objectives. The first is to establish a protocol
for auditing regulatory information. Its object is agreement between a model's
answer and the value published in the official register, for tasks an
environmental agency delegates; it does not measure model capability in general,
reasoning quality, or usefulness for other purposes. The protocol is the
contribution that transfers, and the domain is where we demonstrate it. The
second is to test the mitigations available to an administration. A public body
facing poor answers about its own country can declare the official role in the
prompt, ask in the national language, or run a regionally fine-tuned open
model, and we test all three as designed factors. The regional model we could
run is a community fine-tune of a 7B base on Brazilian Portuguese, the kind an
agency can download; the finding concerns that class, and the released protocol
allows any other candidate to be tested. We do not test retrieval augmentation,
fine-tuning, or human-in-the-loop verification, each of which requires
engineering capacity the prompt-level fixes do not. All queries are answered from
the model's parametric memory, without web search or document retrieval, which
is why the title speaks of recalled rather than retrieved information. The third
objective is to locate where reliability differs and, as far as the design
allows, why; we test the corpus-representation account and claim no causal
identification.

Throughout, the unit of analysis is the model as served, that is, weights,
vendor and serving infrastructure reached through a pinned endpoint, because a
public manager interacts with that stack and any measured gap is a property of
it. Hypotheses, sampling and analysis plan were fixed before data
collection~\citep{nosek2018preregistration} for 15 countries and then extended
to 25; the plan was never deposited publicly, so we do not claim
pre-registration, and the study is reported as exploratory throughout.

\subsection*{Contributions of this work}

\begin{enumerate}\setlength{\itemsep}{2pt}
\item \emph{An auditing protocol for regulatory information.} The question an
  administration faces is whether the value a model returns matches the
  official register. We build five task types from the everyday work of an
  environmental agency (Table~\ref{tab:tasks}), each keyed to an official
  source, so that correctness is adjudicable rather than a matter of expert
  impression, and the protocol transfers to any regulated domain where the
  correct answer is published.
\item \emph{The three localisation fixes agencies reach for, tested rather than
  assumed.} Persona framing, native-language prompting and a regionally tuned
  model are each a procurement or workflow decision. None of them works, and
  they fail differently: asking in the country's own language lowers accuracy
  by $\nnativa$ percentage points, most sharply for the language with the
  scarcest corpus; the role frame is equivalent to no frame within two
  percentage points; and the regionally tuned model recalls the binding value no
  more often than the base model it was tuned from.
\item \emph{A design that separates who is poorly served from why, and reports
  where it cannot.} Countries are stratified on the UNCTAD developed/developing
  split, the Joshi linguistic-resource taxonomy, and development, so that
  geographic, linguistic and economic explanations can be told apart. Between
  countries, coverage and development are collinear; within countries, what
  predicts the country-specificity deficit at the country level is development,
  not the encyclopedic coverage the protocol pre-specified, and the two cannot
  be separated in this sample. We also document that the pre-specified corpus proxy
  assigns the English Wikipedia to every sample country with English as an
  official language, so that it tracks colonial history rather than corpus size.
\item \emph{Ground truth adjudicated by code, and what that changes.} Rebuilding
  two task keys from official registers and moving range comparison into code
  cut the development gradient from $\rho=0.51$ to $\rho=\nrho$ and more than
  doubled the native-language penalty on the same responses; moving the third
  register-valued task from a judge to code removed a confound between judge and
  tier. In this literature the ground-truth key, not the sample size, is the
  binding constraint.
\item \emph{Open resources.} Protocol, prompts, ground-truth registries with
  their official sources, analysis code, and the deviation log are released
  under CC-BY-4.0 and MIT, with a one-command reproduction pipeline.
\end{enumerate}
